% ===================================================================== Chapter 7
\chapter{Meshing}\label{ch:mesh}

\section{CFD Mesh Strategy}
Three meshes with the same topology and systematically different resolution were built so that discretisation error could be estimated
later (Chapter~\ref{ch:mi}). The mesh is a structured butterfly (O-grid) hexahedral mesh generated by a Python script and written directly in
the Fluent mesh format; ANSYS Meshing holds the geometry and named selections in the Workbench project. Scripting was chosen for three
reasons: conformality of the interface is guaranteed by construction (fluid and solid share one node table); every sizing control is an
explicit number under version control; and all three levels are regenerated byte-for-byte by one command.

\section{Structured Hexahedral Mesh}
The cross-section consists of a super-ellipse core filled by transfinite interpolation, a graded fluid annulus and a graded solid annulus;
the section is extruded over uniform axial slabs (Figure~\ref{fig:meshxs}a). The core is a rounded square rather than a circle so that its
four edges map onto the annulus without a singular point on the axis. All three levels are 100\,\% hexahedra (Table~\ref{tab:meshfam}).

%% generated by gen_tables.py - RE-ANALYSIS 2026
\begin{table}[htbp]
\centering
\footnotesize
\caption{CFD mesh family: sizing controls (same topology on all levels)}
\label{tab:meshfam}
\begin{tabular}{lrrr}
\toprule
Control & Coarse & Medium & Fine \\
\midrule
Core divisions per side $N_C$ & 8 & 12 & 18 \\
Circumferential divisions $N_\theta = 4N_C$ & 32 & 48 & 72 \\
Fluid radial layers & 18 & 24 & 32 \\
Solid radial layers & 7 & 10 & 15 \\
Axial divisions & 60 & 90 & 135 \\
Axial cell length [mm] & 10.000 & 6.667 & 4.444 \\
First fluid cell height [\textmu m] & 12.2 & 12.2 & 12.2 \\
Fluid growth rate & 1.3182 & 1.2091 & 1.1385 \\
Solid first cell at the interface [mm] & 0.5 & 0.5 & 0.5 \\
\bottomrule
\end{tabular}
\par\smallskip\parbox{\linewidth}{\footnotesize\raggedright Source: \texttt{05\_Meshing/MESH\_NOTES.md} \S3--\S4. Linear refinement ratio $r \approx 1.46$ between levels.}
\end{table}


\begin{figure}[htbp]
\centering
\begin{subfigure}[t]{0.40\linewidth}\centering\includegraphics[width=\linewidth]{mesh_01_cross_section_full.png}\caption{Cross-section (coarse level)}\end{subfigure}\hfill
\begin{subfigure}[t]{0.58\linewidth}\centering\includegraphics[width=\linewidth]{mesh_03_inflation_near_wall.png}\caption{Near-wall grading at the interface}\end{subfigure}
\caption[Structured O-grid CFD mesh]{Structured O-grid CFD mesh: fluid core, graded fluid annulus and solid wall; the first fluid layers next to the interface are drawn true to scale in (b).}
\label{fig:meshxs}
\src{\provmesh}
\end{figure}

\section{Near-Wall Resolution}
A wall-resolved SST solution requires the first cell centre inside the viscous sublayer. The first-layer height of 12.2~\textmu m was derived
from the Section 2 wall-shear estimate for $y^+ \approx 1$. The converged solution showed that the height holds on the four O-grid axes but
is about 15\,\% smaller at the diagonals (about 10.4~\textmu m), a geometric consequence of the super-ellipse core. The resulting
circumferential $y^+$ pattern is geometric rather than physical (the wall shear varies by only 0.078\,\% around the circumference) and is
harmless because the cells are thinner, not thicker, than designed.

\section{Inflation / First Cell}
The mesh has no separate prism stack; instead the whole fluid annulus is a geometrically graded ring stack solved so that it exactly fills
the annulus from the 12.2~\textmu m first layer. The first-layer height is identical on all three levels, so that the near-wall treatment does
not change with refinement. On the medium mesh, 17 layers lie within 1.56~mm of the wall. The solid starts with a 0.5~mm layer at the interface.

\section{Mesh Quality}
Mesh quality was reported by Fluent (\texttt{mesh/check}, \texttt{mesh/quality}) and reproduced independently by a Python audit to within
$10^{-6}$ relative (Table~\ref{tab:meshq}, Figure~\ref{fig:meshq}). All minima are well above Fluent's orthogonal-quality threshold of 0.1.
The high aspect ratios are intended (thin wall-parallel cells in the near-wall layers). On the fine mesh 540 O-grid corner cells (0.108\,\%)
exceed a skewness of 0.80; the minimum orthogonal quality therefore decreases with refinement, which was located and accepted.

%% generated by gen_tables.py - RE-ANALYSIS 2026
\begin{table}[htbp]
\centering
\footnotesize
\caption{CFD mesh statistics and quality (Fluent \texttt{mesh/check} and \texttt{mesh/quality})}
\label{tab:meshq}
\begin{tabular}{lrrr}
\toprule
Quantity & Coarse & Medium (baseline) & Fine (reference) \\
\midrule
Total cells & 51,840 & 159,840 & 500,580 \\
Fluid / solid cells & 38,400 / 13,440 & 116,640 / 43,200 & 354,780 / 145,800 \\
Conformal interface faces & 1,920 & 4,320 & 9,720 \\
Cell type & 100\,\% hexahedra & 100\,\% hexahedra & 100\,\% hexahedra \\
Minimum orthogonal quality & 0.4482 & 0.4458 & 0.3144 \\
Maximum skewness & 0.6564 & 0.7600 & 0.8364 \\
Maximum aspect ratio & 974.5 & 654.4 & 437.9 \\
Negative-volume cells & 0 & 0 & 0 \\
\bottomrule
\end{tabular}
\par\smallskip\parbox{\linewidth}{\footnotesize\raggedright Source: \texttt{05\_Meshing/MESH\_NOTES.md} \S5--\S6 and \texttt{Mesh\_Quality/MESH\_QUALITY\_TABLE.md}; cell totals \texttt{MASTER\_PROJECT\_DATA.csv} rows M033, M050, M058.}
\end{table}


\section{Coarse / Medium / Fine Meshes}
Refinement is applied in all three directions at once, giving linear refinement ratios of 1.455 and 1.463 between the levels, above the value
of 1.3 recommended for Richardson extrapolation \cite{celik2008}. Figure~\ref{fig:meshref} compares the cross-sections.

\begin{figure}[htbp]
\centering
\includegraphics[width=0.80\linewidth]{mesh_08_refinement_comparison.png}
\caption[Systematic refinement of the CFD mesh family]{Systematic refinement of the CFD mesh family: coarse (51,840~cells), medium (159,840~cells, baseline) and fine (500,580~cells).}
\label{fig:meshref}
\src{\provmesh}
\end{figure}

\section{Conformal Fluid-Solid Interface}
The interface has 1,920 / 4,320 / 9,720 faces on the three levels, each with exactly one fluid owner and one solid neighbour; there are no
orphan faces or hanging nodes, and the interface node radius is exact to $1.7\times10^{-18}$ m (Figure~\ref{fig:meshq}a). On reading the mesh,
Fluent creates the coupled \texttt{fluid\_solid\_interface-shadow} pair automatically, confirming the conformal treatment.

\begin{figure}[htbp]
\centering
\begin{subfigure}[t]{0.56\linewidth}\centering\includegraphics[width=\linewidth]{mesh_04_interface_conformity.png}\caption{Interface conformity audit (coarse)}\end{subfigure}\hfill
\begin{subfigure}[t]{0.42\linewidth}\centering\includegraphics[width=\linewidth]{mesh_10_orthogonal_quality_hist.png}\caption{Orthogonal-quality distribution}\end{subfigure}
\caption[Interface conformity and orthogonal quality of the CFD meshes]{Interface conformity and orthogonal quality of the CFD meshes; dotted lines in (b) mark the minima reported by Fluent.}
\label{fig:meshq}
\src{\provmesh}
\end{figure}

\section{Student-License Constraint}
Fluent Student limited the solver to 4-way parallel execution but stated no cell limit; the 500,580-cell fine mesh was read, checked and
solved. The binding constraint came from the structural side: ANSYS Mechanical Student refuses models above 128,000~nodes (measured:
127,955~nodes solved, 128,045 refused). The fine CFD solid would correspond to about 157\,k structural nodes at 1:1 resolution, beyond that
limit.

\section{Mesh Selection}
The medium mesh (159,840~cells) was selected as the official baseline for all downstream work, and the fine mesh was kept as the
verification reference (decision D-034). The quantitative basis of this choice -- a small and conservative error for the structural phase,
at a third of the cost of the fine mesh -- is given in Chapter~\ref{ch:mi}.
